\chapter{Cluster Implementation and Security Policies}

This chapter documents the practical implementation of the Kubernetes cluster and the policy engine that constitute the security core of the environment designed in Chapter~5. Unlike that chapter, which establishes architectural decisions and their theoretical justification, this chapter focuses on concrete execution: the steps carried out, the deployed policies with their complete manifests, the executed tests and the observed results. The implementation of the DevSecOps pipeline automation components --- ArgoCD for continuous deployment and GitHub Actions for continuous integration --- is addressed in the following chapters.

The chapter is structured into seven sections. The first covers Kubernetes cluster provisioning on the laboratory environment. The second describes Kyverno deployment and the verification of its integration with the API Server. The third, fourth and fifth sections document respectively the implementation of validation, mutation and generation policies, with their manifests, technical justifications and associated tests. The sixth section collects the most relevant operational considerations and empirical findings of the process, including the verification of webhook execution order and the integral \textit{auto-hardening} test. The chapter closes with the seventh section, dedicated to Policy Reports as the native compliance audit mechanism.

% ─────────────────────────────────────────────
\section{Kubernetes cluster provisioning}
% ─────────────────────────────────────────────

The starting point of the project is the provisioning of a functional Kubernetes cluster on the laboratory environment described in Section~4.2. The virtual machines, provided by the supervisor, run Ubuntu Server 24.04~LTS and are pre-configured with the necessary packages (\texttt{kubeadm}, \texttt{kubelet}, \texttt{kubectl}) and with the VirtualBox network layer in NAT plus \textit{host-only adapter} mode.

The control plane is initialised from the \texttt{control} node using the following command:

\begin{lstlisting}
sudo kubeadm init --apiserver-advertise-address=192.168.1.101
\end{lstlisting}

The \texttt{-{}-apiserver-advertise-address} parameter is necessary in environments with multiple network interfaces to indicate to kubeadm on which address the API Server should listen. Without this parameter, kubeadm would select the default interface, which in VirtualBox corresponds to the NAT adapter (\texttt{10.0.2.x}), inaccessible from the worker node. By specifying the \textit{host-only} interface address, it is guaranteed that cluster components --- including Kyverno's webhooks --- are reachable through the laboratory's internal network.

After initialisation, \texttt{kubectl} access is configured for the \texttt{student} user by copying the configuration file generated by kubeadm:

\begin{lstlisting}
mkdir -p $HOME/.kube
sudo cp -i /etc/kubernetes/admin.conf $HOME/.kube/config
sudo chown $(id -u):$(id -g) $HOME/.kube/config
\end{lstlisting}

Next, the worker node (\texttt{192.168.1.102}) joins the cluster via the \texttt{kubeadm join} command generated during initialisation, which includes the bootstrap token and the API Server certificate hash to guarantee connection authenticity.

Once both nodes form part of the cluster, the Calico CNI plugin version 3.28 is installed:

\begin{lstlisting}
kubectl apply -f https://raw.githubusercontent.com/projectcalico/calico/v3.28.0/manifests/calico.yaml
\end{lstlisting}

Calico configures the pod addressing space with the CIDR \texttt{172.16.0.0/16}, the default value for this distribution. Once installation is complete and after system component convergence, the cluster reaches operational state with both nodes in \texttt{Ready} state.

% ─────────────────────────────────────────────
\section{Kyverno deployment and integration}
% ─────────────────────────────────────────────

With the cluster operational, the next step is the installation of Kyverno as the security policy engine. Kyverno is deployed via Helm in the dedicated namespace \texttt{kyverno}:

\begin{lstlisting}
helm repo add kyverno https://kyverno.github.io/kyverno/
helm repo update
helm install kyverno kyverno/kyverno -n kyverno --create-namespace
\end{lstlisting}

After installation, the verification of pod state in the \texttt{kyverno} namespace confirms that the four controllers are active and in \texttt{Running} state, as shown in Figure~\ref{fig:kyverno-pods}.

\begin{figure}[htbp]
\centering
\includegraphics[width=0.85\linewidth]{kyverno/compKyverno.png}
\caption{Kyverno controllers deployed in the cluster}
\label{fig:kyverno-pods}
\end{figure}

The four controllers --- \texttt{admission-controller}, \texttt{background-controller}, \texttt{cleanup-\allowbreak controller} and \texttt{reports-controller} --- operate independently and complementarily, as described in Section~5.2.2. The most relevant controller for the system's real-time behaviour is the \textit{admission controller}, which registers in the API Server immediately after its deployment.

The verification of registered webhooks confirms that Kyverno has correctly established its integration with the API Server. The query of \texttt{Validating\allowbreak{}WebhookConfiguration} objects shows seven configurations registered by Kyverno, as seen in Figure~\ref{fig:kyverno-webhooks}.

\begin{figure}[htbp]
\centering
\includegraphics[width=0.85\linewidth]{kyverno/webhooksKyverno.png}
\caption{ValidatingWebhookConfigurations registered by Kyverno in the API Server}
\label{fig:kyverno-webhooks}
\end{figure}

Among the registered configurations, the most relevant for this work is \texttt{kyverno-resource-\allowbreak{}validating-\allowbreak{}webhook-\allowbreak{}cfg}, which intercepts creation and modification requests for arbitrary cluster resources and routes them to Kyverno's \textit{admission controller} for evaluation. The additional configurations cover specific cases such as policy validation itself (\texttt{kyverno-policy-\allowbreak{}validating-\allowbreak{}webhook-\allowbreak{}cfg}), global context (\texttt{kyverno-global-context-\allowbreak{}validating-\allowbreak{}webhook-\allowbreak{}cfg}) and exception and TTL management. Similarly, Kyverno registers the corresponding \texttt{Mutating\allowbreak{}Webhook\allowbreak{}Configuration} objects for processing \textit{mutate} policies.

With Kyverno installed and its webhooks active, the cluster is ready to receive security policies. Figure~\ref{fig:estado-inicial} shows the cluster state at this point, before applying any policy: the \texttt{ClusterPolicy} query returns no resources.

\begin{figure}[htbp]
\centering
\includegraphics[width=0.58\linewidth]{kyverno/estadoInicial.png}
\caption{Initial cluster state: no policies deployed}
\label{fig:estado-inicial}
\end{figure}

% ─────────────────────────────────────────────
\subsection{Kyverno's autogen mechanism}
\label{subsec:autogen}
% ─────────────────────────────────────────────

Before documenting the implemented policies, it is necessary to understand a mechanism that conditions the real scope of all of them: automatic rule generation, known as \textit{autogen}.

\subsubsection*{The problem: Pods are rarely created directly}

The eight policies in this work declare \texttt{match.resources.kinds:\,[Pod]}, which appears to restrict their application to \texttt{Pod} type resources. However, in practice Pods are not created directly: they are generated by higher-level controllers --- \texttt{Deployment}, \texttt{StatefulSet}, \texttt{DaemonSet}, \texttt{Job}, \texttt{CronJob}, \texttt{ReplicaSet} --- through a \texttt{spec.template} field that contains the specification of the Pod to be instantiated.

If Kyverno only evaluated resources of type \texttt{Pod}, a developer --- or an attacker --- could create a \texttt{Deployment} with insecure configuration without the policy intercepting it. The Pods derived from that Deployment would be individually rejected by the \textit{admission webhook} when the controller tried to create them, but the \texttt{Deployment} object would remain persisted in etcd and the controller would enter a loop of recurring errors, generating operational noise and leaving the root resource in an inconsistent state. In permissive failure modes, depending on the webhook's \texttt{failurePolicy} configuration, some resources might even manage to execute.

\subsubsection*{The solution: auto-generated rules}

Kyverno solves this problem transparently for the policy author. When processing a \texttt{ClusterPolicy} that declares \texttt{match.kinds:\,[Pod]}, the engine analyses the defined pattern and automatically generates additional rules --- called \textit{autogen rules} --- that apply the same control over the \texttt{spec.template.spec} of controllers that create Pods. The developer defines the policy once and Kyverno internally replicates it to cover \texttt{Deployment}, \texttt{StatefulSet}, \texttt{DaemonSet}, \texttt{Job}, \texttt{CronJob} and \texttt{ReplicaSet}.

Generation distinguishes two cases depending on nesting depth. For most controllers, the auto-generated rule is called \texttt{autogen-<original-name>} and applies the pattern over \texttt{spec.template.spec}. For \texttt{CronJob}, whose Pod specification is one level deeper (\texttt{spec.jobTemplate.\allowbreak{}template.spec}), Kyverno generates a second additional rule, \texttt{autogen-cronjob-<original-name>}, which adjusts the pattern path accordingly.

The empirical verification of this behaviour is carried out by inspecting the \texttt{ClusterPolicy} object once deployed:

\begin{lstlisting}
kubectl get clusterpolicy disallow-root -o yaml | grep -A 50 "autogen:"
\end{lstlisting}

The output of this command shows the rules \texttt{autogen-require-run-as-non-root}, covering Deployments, StatefulSets, Jobs and CronJobs, and \texttt{autogen-cronjob-require-run-as-non-root}, handling the specific \texttt{CronJob} case. These rules are generated internally by Kyverno and stored in the policy object in etcd; they do not require any manual intervention from the operator.

\subsubsection*{Differentiated behaviour: validate vs. mutate}

The autogen mechanism does not behave identically for all policy types:

\begin{itemize}
  \item For \textbf{validate} policies, autogen is always automatic and unconditional. Kyverno generates additional rules for all controllers that create Pods without the need for any explicit configuration.
  \item For \textbf{mutate} policies using \texttt{patchStrategicMerge} --- such as the four mutation policies in this work (P4, P5, P6, P7) --- autogen is also automatic. Kyverno adapts the patch so that it applies over the controller's \texttt{spec.template.spec} field, guaranteeing that mutations injected on direct Pods are equally replicated on Pods derived from Deployments or StatefulSets.
  \item For \textbf{mutate} policies using \texttt{patchesJson6902}, autogen is not applied by default, since the absolute JSON paths of this mechanism are not automatically generalisable. In these cases, the operator must explicitly configure the autogen behaviour.
\end{itemize}

\subsubsection*{Implication for this work's policies}

The eight implemented policies --- P1 to P8 --- declare \texttt{match.kinds:\,[Pod]}, but thanks to the autogen mechanism they effectively cover Deployments, StatefulSets, Jobs and CronJobs. This means that the real scope of the security system is significantly greater than what can be derived from a direct reading of the manifests.

This behaviour became evident during the initial implementation phases: PolicyViolations generated by the Reports Controller included messages such as \texttt{autogen-require-run-as-non-root} being applied to \texttt{Deployment} and \texttt{ReplicaSet} objects of Kyverno's own components and \texttt{kube-system}, confirming that coverage automatically extended to higher-level resources without having declared those kinds in the policy.

Kyverno allows controlling this behaviour through the annotation \texttt{pod-policies.kyverno.io/\allowbreak{}autogen-controllers} in the policy's \texttt{metadata}, which accepts an explicit list of controllers to cover or the value \texttt{none} to completely disable autogen. In this work the default behaviour is used, which offers the broadest possible coverage without additional configuration.

% ─────────────────────────────────────────────
\section{Validation policies}
% ─────────────────────────────────────────────

Validation policies constitute the reactive layer of the security system: they act by rejecting resources that violate conditions that cannot be resolved through automatic mutation. As justified in Section~5.5.4, mutation policies have precedence and act first; validation policies therefore operate as a final barrier for explicitly insecure configurations that should not be silently corrected, but rejected with a clear message to the requester.

\subsection{Policy P1: \texttt{disallow-root}}

The first validation policy prevents the deployment of pods that do not explicitly specify that their containers must run as a non-root user:

\begin{lstlisting}
apiVersion: kyverno.io/v1
kind: ClusterPolicy
metadata:
  name: disallow-root
spec:
  validationFailureAction: Enforce
  background: true
  rules:
  - name: require-run-as-non-root
    match:
      any:
      - resources:
          kinds:
          - Pod
    validate:
      message: "Containers must not run as root"
      pattern:
        spec:
          securityContext:
            runAsNonRoot: true
\end{lstlisting}

The policy operates through \textit{pattern matching}: Kyverno verifies that the field \texttt{spec.security\allowbreak{}Context.\allowbreak{}runAsNonRoot} exists with the value \texttt{true}. If the field is absent or is \texttt{false}, the webhook returns a validation error with the defined message and the request is rejected.

The technical justification for this control is twofold. On one hand, a process with \texttt{UID~0} inside the container inherits the complete set of Linux capabilities assigned to it, including \texttt{CAP\_SYS\_ADMIN} and \texttt{CAP\_NET\_RAW}, which significantly expand the attack surface. On the other hand, if a container isolation escape occurs --- through a runtime or kernel vulnerability --- a process that was already running as root in the container will obtain root privileges on the host node, potentially compromising the entire cluster.

The validation test applies the following pod, which specifies no \texttt{securityContext}:

\begin{lstlisting}
apiVersion: v1
kind: Pod
metadata:
  name: nginx-insecure
spec:
  containers:
  - name: nginx
    image: nginx
\end{lstlisting}

The result, shown in Figure~\ref{fig:test1}, confirms that Kyverno's webhook rejects the request indicating the exact path of the failing field.

\begin{figure}[htbp]
\centering
\includegraphics[width=0.95\linewidth]{kyverno/test1.png}
\caption{P1: pod without \texttt{runAsNonRoot} blocked by Kyverno's webhook}
\label{fig:test1}
\end{figure}

\subsection{Policy P2: \texttt{disallow-privilege-escalation}}

The second policy requires all containers to explicitly declare \texttt{allowPrivilege\allowbreak{}Escalation: false} in their container-level \texttt{securityContext}:

\begin{lstlisting}
apiVersion: kyverno.io/v1
kind: ClusterPolicy
metadata:
  name: disallow-privilege-escalation
spec:
  validationFailureAction: Enforce
  rules:
  - name: no-privilege-escalation
    match:
      any:
      - resources:
          kinds:
          - Pod
    validate:
      message: "Privilege escalation is not allowed"
      pattern:
        spec:
          containers:
          - securityContext:
              allowPrivilegeEscalation: false
\end{lstlisting}

This field controls the \texttt{no\_new\_privs} flag of the Linux kernel \cite{linux_no_new_privs}, which prevents any child process from acquiring more privileges than its parent process. Without this flag active, an unprivileged \texttt{UID} process could invoke a binary with the \texttt{SUID} bit active --- such as \texttt{sudo} --- to elevate its privileges to root, circumventing \texttt{runAsNonRoot} controls.

The test applies a pod that explicitly requests privilege escalation. The result, in Figure~\ref{fig:test2}, shows rejection with the path of the offending field.

\begin{figure}[htbp]
\centering
\includegraphics[width=0.95\linewidth]{kyverno/test2.png}
\caption{P2: pod with \texttt{allowPrivilegeEscalation: true} blocked}
\label{fig:test2}
\end{figure}

It is worth noting the interaction between this policy and mutation policy P5. Since \textit{mutating webhooks} are invoked before \textit{validating webhooks}, P5 automatically injects \texttt{allow\-Privilege\-Escalation:\,false} in pods that do not specify this field. Consequently, P2 acts only as a final barrier for cases where the value \texttt{true} has been explicitly declared --- an insecure intention declaration that should not be silently corrected, but rejected.

\subsection{Policy P3: \texttt{restrict-image-registries}}

The third policy restricts image registries to a set of verified sources, implemented via Kyverno's \texttt{foreach} mechanism:

\begin{lstlisting}
apiVersion: kyverno.io/v1
kind: ClusterPolicy
metadata:
  name: restrict-image-registries
spec:
  validationFailureAction: Enforce
  background: true
  rules:
  - name: validate-registries
    match:
      any:
      - resources:
          kinds:
          - Pod
    validate:
      message: >-
        Only images from authorised registries are allowed:
        docker.io, ghcr.io, registry.k8s.io, quay.io.
        The image {{ element.image }} is not permitted.
      foreach:
      - list: "request.object.spec.containers"
        deny:
          conditions:
            all:
            - key: "{{ element.image }}"
              operator: AnyNotIn
              value:
              - "docker.io/*"
              - "ghcr.io/*"
              - "registry.k8s.io/*"
              - "quay.io/*"
              - "nginx*"
              - "alpine*"
              - "busybox*"
              - "nginxinc/*"
\end{lstlisting}

The policy iterates over all pod containers and denies the request if any image does not belong to the authorised patterns. Patterns without registry prefix (\texttt{nginx*}, \texttt{alpine*}, \texttt{busybox*}) are necessary because Kubernetes implicitly resolves unprefixed references to Docker Hub. The inclusion of \texttt{ghcr.io/*} is an operational requirement, given that Kyverno and other cluster component images are distributed through that registry.

This policy constitutes a first line of defence against supply chain attacks of the type documented by the TeamPCP group \cite{teampcp2026}: an attacker who hosts a malicious image on an arbitrary registry will not be able to run it in the cluster without modifying this policy. In mature production environments, the control would be complemented with digital signature verification using Cosign or Notary, considered future work in Section~4.7.

During implementation, a relevant engine behaviour was observed: the variable \texttt{\{\{ element.image \}\}} referenced in the top-level \texttt{message} field does not expand at rejection time, since that variable only has scope within the \texttt{foreach} block. Kyverno accepts the policy as valid and the \texttt{READY} field shows \texttt{True}, but at rejection time the error message only shows the static part, omitting the specific image name. This behaviour illustrates Kyverno's variable scoping model and must be taken into account when designing informative error messages in policies with \texttt{foreach}.

The validation test applies a pod with an image from an unauthorised registry. The result, in Figure~\ref{fig:test3}, reveals a defence-in-depth behaviour: the pod is simultaneously blocked by two independent policies.

\begin{figure}[htbp]
\centering
\includegraphics[width=0.95\linewidth]{kyverno/test3.png}
\caption{P3: simultaneous blocking by \texttt{disallow-root} and \texttt{restrict-image-registries}}
\label{fig:test3}
\end{figure}

This result demonstrates that the policy system acts cumulatively: the same resource can violate multiple policies simultaneously, and the error message returned enumerates all detected violations. The image \texttt{evilregistry.io/backdoor:latest} is blocked both for not specifying \texttt{runAsNonRoot} (P1) and for coming from an unauthorised registry (P3), two independent risk vectors detected in a single response.

% ─────────────────────────────────────────────
\section{Mutation policies}
% ─────────────────────────────────────────────

Mutation policies constitute the main security mechanism of the system, aligning with the central hypothesis of the work: that the automatic correction of insecure configurations at admission time offers greater resilience and better development experience than mere rejection. Instead of blocking pods with incomplete configurations, Kyverno automatically enriches them with the correct security profile, making \textit{hardening} transparent to the developer.

The \texttt{patchStrategicMerge} mechanism, used by the four policies in this block, performs an intelligent merge of the original object with the defined patch: it adds absent fields and respects values already explicitly declared by the developer, guaranteeing that mutations are additive and non-destructive.

\subsection{Policy P4: \texttt{add-default-resources}}

The fourth policy automatically injects CPU and memory \texttt{requests} and \texttt{limits} in all containers that do not declare them:

\begin{lstlisting}
apiVersion: kyverno.io/v1
kind: ClusterPolicy
metadata:
  name: add-default-resources
spec:
  rules:
  - name: add-resources
    match:
      any:
      - resources:
          kinds:
          - Pod
    mutate:
      patchStrategicMerge:
        spec:
          containers:
          - (name): "*"
            resources:
              requests:
                cpu: "100m"
                memory: "128Mi"
              limits:
                cpu: "500m"
                memory: "256Mi"
\end{lstlisting}

The simultaneous injection of \texttt{requests} and \texttt{limits} with the same values --- rather than only \texttt{limits} --- responds to a technically justified decision. Kubernetes classifies pods into three Quality of Service classes (\textit{QoS}) \cite{kubernetes_qos}: \texttt{Guaranteed} (when \texttt{requests\,==\,limits} for all containers), \texttt{Burstable} and \texttt{BestEffort}. A pod with \texttt{Guaranteed} QoS is the last candidate for eviction when the node experiences resource pressure. Additionally, without \texttt{requests} defined, the Scheduler lacks information for \textit{bin-packing} decisions and may schedule pods on nodes without sufficient capacity, resulting in OOM failures at runtime.

The test deploys a pod without any \texttt{resources} field defined. As shown in Figure~\ref{fig:test4}, Kyverno automatically injects \texttt{requests} and \texttt{limits} values, visible in the \texttt{kubectl get pod -o yaml} output.

\begin{figure}[htbp]
\centering
\includegraphics[width=0.95\linewidth]{kyverno/test4.png}
\caption{P4: automatic injection of \texttt{resources.requests} and \texttt{resources.limits}}
\label{fig:test4}
\end{figure}

It is especially significant that the output also shows the labels \texttt{security.tfm/\allowbreak{}hardened-by:\,kyverno} and \texttt{security.tfm/\allowbreak{}policy-version:\,v1.0}, evidence that P7 has also acted on the same resource in the same admission cycle. This confirms that multiple mutation policies are applied in a coordinated manner to each object in a single admission pass.

\subsection{Policy P5: \texttt{add-default-securitycontext}}

The fifth policy injects a secure \texttt{securityContext} in all containers that do not specify one at the container level:

\begin{lstlisting}
apiVersion: kyverno.io/v1
kind: ClusterPolicy
metadata:
  name: add-default-securitycontext
spec:
  rules:
  - name: add-container-security-context
    match:
      any:
      - resources:
          kinds:
          - Pod
    mutate:
      patchStrategicMerge:
        spec:
          containers:
          - (name): "*"
            securityContext:
              allowPrivilegeEscalation: false
              capabilities:
                drop:
                - ALL
\end{lstlisting}

This policy implements two critical \textit{hardening} controls automatically. The first, \texttt{allow\-Privilege\-Escalation:\,false}, activates the kernel's \texttt{no\_new\_privs} flag, preventing privilege escalation through SUID binaries. Here the policy acts \textit{proactively} --- injecting the secure configuration before reaching validation --- in contrast to P2, which acts \textit{reactively} by rejecting the explicit declaration of the insecure value.

The second control, \texttt{capabilities.drop:\,ALL}, removes the complete set of Linux capabilities inherited by default. Without additional configuration, a container inherits capabilities such as \texttt{NET\_RAW} (raw packet construction), \texttt{MKNOD} (block device creation) or \texttt{CHOWN} (file ownership change), which unnecessarily expand the attack surface. By dropping them all by default, developers are forced to explicitly declare the capabilities their applications actually need.

This policy is the most representative example of the \textit{auto-hardening} concept: the developer can deploy a pod with minimal configuration without deep knowledge of Kubernetes security fields, and the platform guarantees a solid security profile in a completely transparent manner.

The test deploys a pod with \texttt{securityContext} defined only at the pod level. Figure~\ref{fig:test5} shows the result: the container-level \texttt{securityContext} contains \texttt{allow\-Privilege\-Escalation:\,false}, \texttt{capabilities.drop:\,[ALL]} and \texttt{readOnly\-Root\-Filesystem:\,true} --- the latter injected by P6 in the same admission cycle.

\begin{figure}[htbp]
\centering
\includegraphics[width=0.95\linewidth]{kyverno/test5.png}
\caption{P5: automatic injection of \texttt{allowPrivilegeEscalation} and \texttt{capabilities.drop}}
\label{fig:test5}
\end{figure}

\subsection{Policy P6: \texttt{add-readonly-rootfs}}

The sixth policy forces the root filesystem of all containers to read-only mode:

\begin{lstlisting}
apiVersion: kyverno.io/v1
kind: ClusterPolicy
metadata:
  name: add-readonly-rootfs
spec:
  rules:
  - name: add-readonly-filesystem
    match:
      any:
      - resources:
          kinds:
          - Pod
    mutate:
      patchStrategicMerge:
        spec:
          containers:
          - (name): "*"
            securityContext:
              readOnlyRootFilesystem: true
\end{lstlisting}

Most post-exploitation techniques after compromising a container --- downloading reverse shells, installing cryptominers, executing reconnaissance tools --- require writing files to the container filesystem. A read-only root filesystem makes these attempts fail immediately, containing the blast radius of the intrusion.

The test verifies this effect in two steps, shown in Figure~\ref{fig:test6}. First, the \texttt{kubectl get pod -o yaml} output confirms that the \texttt{readOnly\-Root\-Filesystem:\,true} field has been injected. Second, a command is executed inside the container that attempts to create a file in the root directory:

\begin{lstlisting}
kubectl exec test-mutate-seccontext -- touch /test-file
touch: /test-file: Read-only file system
command terminated with exit code 1
\end{lstlisting}

\begin{figure}[htbp]
\centering
\includegraphics[width=0.90\linewidth]{kyverno/test6.png}
\caption{P6: verification of read-only filesystem at runtime}
\label{fig:test6}
\end{figure}

This result is especially valuable: it demonstrates that the mutation applied by Kyverno at admission time has a real and verifiable effect on container behaviour at runtime, not just on the manifest stored in etcd.

A known operational limitation is that certain applications require writing to directories such as \texttt{/tmp} or \texttt{/var/cache}. The correct solution --- applied in the integral test of Section~6.6.2 --- consists of mounting an \texttt{emptyDir} volume at the specific paths that require write access. This practice does not weaken the policy, but forces developers to be explicit about their write requirements, which is itself a good container design practice.

\subsection{Policy P7: \texttt{add-default-labels}}

The seventh policy injects traceability labels in all pods processed by the system:

\begin{lstlisting}
apiVersion: kyverno.io/v1
kind: ClusterPolicy
metadata:
  name: add-default-labels
spec:
  rules:
  - name: add-security-labels
    match:
      any:
      - resources:
          kinds:
          - Pod
    mutate:
      patchStrategicMerge:
        metadata:
          labels:
            security.tfm/hardened-by: kyverno
            security.tfm/policy-version: "v1.0"
\end{lstlisting}

Although less obvious than direct security policies, this policy implements an observability capability relevant in production environments. The injected labels allow identifying at any time which pods have been processed by the system and with which policy version, facilitating coverage auditing. The \texttt{security.tfm/\allowbreak{}policy-version} field acts as a marker that allows correlating the applied security profile with the policy version in force at the time of deployment.

Verification via \texttt{kubectl get pods --show-labels}, shown in Figure~\ref{fig:test7}, confirms that the labels are present in all active pods in the \texttt{default} namespace.

\begin{figure}[htbp]
\centering
\includegraphics[width=0.90\linewidth]{kyverno/test7.png}
\caption{P7: traceability labels automatically injected in pods}
\label{fig:test7}
\end{figure}

% ─────────────────────────────────────────────
\section{Generation policy}
% ─────────────────────────────────────────────

\subsection{Policy P8: \texttt{generate-networkpolicy-metadata}}

The eighth policy demonstrates the third type of action available in Kyverno --- \textit{generate} --- with no functional equivalent in OPA Gatekeeper. Unlike previous policies, which act on resources in the process of creation, this one reacts to the creation of a \texttt{Namespace} and automatically generates a derived resource: a NetworkPolicy that blocks access to the cloud instance metadata service (IMDS):

\begin{lstlisting}
apiVersion: kyverno.io/v1
kind: ClusterPolicy
metadata:
  name: generate-networkpolicy-metadata
spec:
  rules:
  - name: block-metadata-service
    match:
      any:
      - resources:
          kinds:
          - Namespace
    generate:
      apiVersion: networking.k8s.io/v1
      kind: NetworkPolicy
      name: block-cloud-metadata
      namespace: "{{request.object.metadata.name}}"
      synchronize: true
      data:
        spec:
          podSelector: {}
          policyTypes:
          - Egress
          egress:
          - to:
            - ipBlock:
                cidr: 0.0.0.0/0
                except:
                - 169.254.169.254/32
\end{lstlisting}

The generated NetworkPolicy allows all outgoing traffic (\texttt{cidr:\,0.0.0.0/0}) except towards the address \texttt{169.254.169.254/32}. This IP corresponds to the IMDS present in the main cloud providers (AWS, GCP, Azure), accessible from any instance without prior authentication. An attacker who compromises a pod can make a request to that endpoint to obtain temporary IAM credentials, service tokens and underlying instance network configuration, enabling escalation from the pod to the cloud account. This vector was documentedly used in the Capital One incident in 2019, where SSRF was exploited to access the IMDS and exfiltrate sensitive data \cite{capitalone2019}.

The \texttt{synchronize:\,true} field activates the Background Controller's \textit{drift protection} mechanism: if the generated NetworkPolicy is modified or deleted, Kyverno automatically restores it.

The first test, shown in Figure~\ref{fig:test8}, demonstrates automatic generation: namespace \texttt{test-networkpolicy} is created and just 9~seconds later the NetworkPolicy \texttt{block-cloud-metadata} exists in that namespace without having been manually applied. The detailed description of the generated resource confirms that Kyverno has added the corresponding traceability labels and that the egress rule exclusively blocks IP \texttt{169.254.169.254/32}.

\begin{figure}[htbp]
\centering
\includegraphics[width=0.95\linewidth]{kyverno/test8.png}
\caption{P8: NetworkPolicy automatically generated upon namespace creation}
\label{fig:test8}
\end{figure}

The second test, shown in Figure~\ref{fig:test9}, demonstrates the \textit{drift protection} mechanism: the NetworkPolicy is manually deleted via \texttt{kubectl delete} and, 57~seconds later, it has been automatically regenerated by the Background Controller, without any human intervention.

\begin{figure}[htbp]
\centering
\includegraphics[width=0.90\linewidth]{kyverno/test9.png}
\caption{P8: \textit{drift protection} --- NetworkPolicy regenerated 57~s after manual deletion}
\label{fig:test9}
\end{figure}

This behaviour illustrates a fundamental difference compared to OPA Gatekeeper, which lacks native support for resource generation and cannot implement this pattern of continuous protection.

% ─────────────────────────────────────────────
\section{Operational considerations and lessons learned}
% ─────────────────────────────────────────────

This section collects the most relevant operational findings that emerged during implementation and that do not fit in the individual documentation of each policy. Some of them are directly related to the autogen mechanism described in Section~\ref{subsec:autogen}: the automatic extension of policies to Deployments and ReplicaSets generated friction with ArgoCD during the GitOps pipeline integration phases, as the policies evaluated resources managed by the ArgoCD controller before they were fully compliant. This behaviour is documented in detail in Chapter~7, in the context of the complete pipeline integration.

\subsection{Exclusion of infrastructure namespaces}

During the initial implementation phases, validation policies generated violations on the cluster's own components. In particular, policy P1 (\texttt{disallow-root}) marked certain pods in the \texttt{kube-system} namespace as non-compliant, as they run as root due to functional requirements. In parallel, the first version of P3 (\texttt{restrict-image-\allowbreak{}registries}) did not include the \texttt{ghcr.io/*} pattern, which generated PolicyViolations on Kyverno's own controllers, whose images are distributed via \texttt{ghcr.io/kyverno/}.

The applied solution was twofold: \texttt{ghcr.io/*} was added to the list of authorised registries in P3, correcting the problem structurally; and the policies were revised to exclude infrastructure namespaces in cases where their application is not pertinent. This is a standard practice in Kyverno production deployments, where security policies must be designed with awareness of the system components that execute them. Figure~\ref{fig:rules-aplicadas} shows the final cluster state with all eight policies in \texttt{Ready} state.

\begin{figure}[H]
\centering
\includegraphics[width=0.95\linewidth]{kyverno/rulesAplicadas.png}
\caption{All eight ClusterPolicies in \texttt{Ready} state after complete deployment}
\label{fig:rules-aplicadas}
\end{figure}

\subsection{Integral auto-hardening test}

The most representative test of the work consists of deploying a pod with deliberately minimal configuration and verifying that the complete set of mutation policies transforms it into a pod with a complete security profile. The initial manifest, shown in Figure~\ref{fig:pod-seguro-yaml}, contains only the fields strictly necessary for operation: image, execution user and an \texttt{emptyDir} volume mounted at \texttt{/tmp} to satisfy policy P6's read-only filesystem requirement.

\begin{figure}[H]
\centering
\includegraphics[width=0.65\linewidth]{kyverno/pod-seguro-completo.png}
\caption{Minimal pod manifest used in the integral test}
\label{fig:pod-seguro-yaml}
\end{figure}

\vspace{-\parskip}

After deployment, pod inspection via \texttt{kubectl get pod pod-seguro-completo -o yaml} returns the complete object as persisted in etcd after passing through Kyverno's webhooks. Figures~\ref{fig:test10-1} and~\ref{fig:test10-2} show the result.

\begin{figure}[H]
\centering
\includegraphics[width=0.95\linewidth]{kyverno/test10-1.png}
\caption{Integral test (part 1): pod metadata and resources after admission}
\label{fig:test10-1}
\end{figure}

\vspace{-\parskip}

\begin{figure}[H]
\centering
\includegraphics[width=0.85\linewidth]{kyverno/test10-2.png}
\caption{Integral test (part 2): complete security profile injected by Kyverno}
\label{fig:test10-2}
\end{figure}


The analysis of the persisted object confirms the presence of all fields injected by the four active mutation policies: \texttt{resources.requests} and \texttt{resources.limits} (P4); \texttt{allow\-Privilege\-Escalation:\,false} and \texttt{capabilities.drop:\,[ALL]} (P5); \texttt{readOnly\-Root\-Filesystem:\,true} (P6); and labels \texttt{security.tfm/\allowbreak{}hardened-by:\,kyverno} and \texttt{security.tfm/\allowbreak{}policy-version:\,v1.0} (P7). Calico annotations additionally confirm that the pod has been assigned IP \texttt{172.16.171.112} in CIDR \texttt{172.16.0.0/16} and scheduled on the \texttt{worker} node.

This test constitutes the direct empirical evidence of the work's central thesis: a developer can define a pod with the minimum fields necessary for its function and obtain, in a completely transparent manner, a correctly secured pod that complies with the security profile defined by the organisation.

\subsection{The mutation-validation order verified empirically}

During testing, a behaviour was observed that empirically illustrates the webhook execution order described theoretically in Section~5.3.2. When attempting to deploy a pod with \texttt{allow\-Privilege\-Escalation:\,true} declared explicitly while both P5 and P2 were simultaneously active, the pod was created successfully despite P2 operating in \texttt{Enforce} mode.

The cause is precise: P5 (\textit{mutating webhook}) executed before P2 (\textit{validating webhook}) and overwrote the \texttt{true} value declared in the manifest with \texttt{false}, so that when P2 evaluated the resource, it already contained the compliant value. The \texttt{patch\-Strategic\-Merge} of P5 overwrites the field even when the developer has explicitly declared it with an insecure value.

This behaviour was verified by temporarily disabling P5 and applying the same pod: in the absence of mutation, P2 correctly rejected the request. Upon restoring P5, the pod was again admitted with the mutated value. This experiment empirically confirms that, in the adopted design, mutation has \textit{absolute precedence} over validation, and that validation policies complement but do not replace mutation policies as the primary security mechanism.

% ─────────────────────────────────────────────
\section{Policy Reports: compliance auditing}
% ─────────────────────────────────────────────

Kyverno's Reports Controller automatically generates \texttt{PolicyReport} resources (namespace-scoped) and \texttt{ClusterPolicyReport} resources (cluster-scoped) that consolidate the results of evaluating all policies over all system resources. These reports constitute the native compliance audit mechanism and allow centrally answering the question: which cluster resources comply with the security policies?

Figure~\ref{fig:clusterpolicyreports} shows the \texttt{ClusterPolicyReport} generated for the \texttt{test-networkpolicy} namespace: the \texttt{Namespace} type resource has passed the P8 evaluation, with result 1~PASS and 0~FAIL.

\begin{figure}[htbp]
\centering
\includegraphics[width=0.95\linewidth]{kyverno/kubectl get clusterpolicyreports.png}
\caption{ClusterPolicyReport: evaluation of the generation policy on the test namespace}
\label{fig:clusterpolicyreports}
\end{figure}

Figure~\ref{fig:policyreports} shows the \texttt{PolicyReport} objects for the \texttt{default} namespace, one per active pod. The three pods --- \texttt{pod-seguro-completo}, \texttt{test-mutate-resources} and \texttt{test-mutate-sec\-context} --- record 7~PASS and 0~FAIL.

\begin{figure}[htbp]
\centering
\includegraphics[width=0.95\linewidth]{kyverno/kubectl get policyreports -A.png}
\caption{PolicyReports of the \texttt{default} namespace: 7 PASS per pod, 0 violations}
\label{fig:policyreports}
\end{figure}

The value 7 corresponds exactly to the number of policies that evaluate \texttt{Pod} type resources (P1 to P7). Policy P8 does not appear in these reports because its scope covers \texttt{Namespace} type resources, and its results are reflected in the \texttt{ClusterPolicyReport} objects of Figure~\ref{fig:clusterpolicyreports}. These reports demonstrate that the implemented system not only acts at admission time, but also provides continuous visibility into the cluster's compliance state --- a fundamental requirement in DevSecOps environments oriented towards auditing and regulatory compliance.
